\documentclass{article}
\usepackage[UTF8]{ctex}
\usepackage{amsmath}
\usepackage{amssymb}
\usepackage{graphicx}
\usepackage{geometry}
\geometry{a4paper, margin=2cm}

\newcommand{\problemname}[1]{\begin{center}\Large\bfseries #1\end{center}}
\usepackage{enumitem}

\begin{document}
\problemname{捡贝壳}

\subsection*{题目描述}
蒜头君去海边玩耍看到了很多贝壳，他想捡一些贝壳装到箱子带回家。他的箱子最多能装 20 斤的东西。现在他已经装了两个贝壳到箱子里，每个贝壳重 a 斤。他的箱子还能再装多少斤的贝壳呢？

\subsection*{输入格式}
输入只有一个实数，对应题目中每枚贝壳的重量 a（斤）。数据范围：0 < 2 × a ≤ 20，即 0 < a ≤ 10。

\subsection*{输出格式}
输出只有一个小数，表示蒜头君的箱子还可以装的贝壳重量（单位：斤），即 20 − 2 × a。

\subsection*{样例输入}
\begin{verbatim}
2
\end{verbatim}

\subsection*{样例输出}
\begin{verbatim}
16
\end{verbatim}


\end{document}
